%%%%%%%%%%%%%%%%%%%%%%%%%%%%%%%%%%%%%%%%%%%%%%%%%%%%%%%%%%%%%%%%%
% 							SETTINGS 							%
%%%%%%%%%%%%%%%%%%%%%%%%%%%%%%%%%%%%%%%%%%%%%%%%%%%%%%%%%%%%%%%%%

\documentclass[12pt]{article}

\usepackage{
	amsmath,
	amssymb,
	geometry
}

\geometry{margin=1in}

\author{Álvaro Fernández Barrero and Vicente Gabriel Piñar Ojeda}
\title{Quaternions}
\date{04-10-2026}

%%%%%%%%%%%%%%%%%%%%%%%%%%%%%%%%%%%%%%%%%%%%%%%%%%%%%%%%%%%%%%%%%
% 							DOCUMENT 							%
%%%%%%%%%%%%%%%%%%%%%%%%%%%%%%%%%%%%%%%%%%%%%%%%%%%%%%%%%%%%%%%%%

\begin{document}
	\maketitle
	
	%%%%%%%%%%%%%%%%%%%%%%%%%%%%%%%%%%%%%%%%%%%%%%%%%%%%%%%%%%%%%%%%%%%%%
	% 							INTRODUCTION 							%
	%%%%%%%%%%%%%%%%%%%%%%%%%%%%%%%%%%%%%%%%%%%%%%%%%%%%%%%%%%%%%%%%%%%%%
	
	\section{Theoretical introduction}
	
	A quaternion $q$ is essentially a number such as $q \in \mathbb{H}$, where:
	
	\[
		\mathbb{H} := \{w + xi + yj + zk : w,x,y,z \in \mathbb{R}, i^2 = j^2 = k^2 = ijk = -1 \}
	\]
	
	%%%%%%%%%%%%%%%%%%%%%%%%%%%%%%%%%%%%%%%%%%%%%%%%%%%%%%%%%%%%%%%%%
	% 							OPERATIONS 							%
	%%%%%%%%%%%%%%%%%%%%%%%%%%%%%%%%%%%%%%%%%%%%%%%%%%%%%%%%%%%%%%%%%
	
	\section{Hamilton's equations}

	\section{Quaternions arithmetic}
	
	As a quaternion is nothing but a number, operations among quaternions of course are allowed. For a quaternion $q_0 = a_0 + b_0i + c_0j + d_0k$ and another quaternion $q_1 = a_1 + b_1i + c_1j + d_1k$:
	
	\[
		\begin{aligned}
			& \lambda q
			= \lambda (w + xi + yj + zk) = \lambda w + \lambda xi + \lambda yj + \lambda zk,
			\forall \lambda \in \mathbb{F}, q \in \mathbb{H}
			\\
			& q_0 + q_1
			= (w_0 + x_0i + y_0j + z_0k) + (w_1 + x_1i + y_1j + z_1k)
			\\
			& \hspace{3em} = w_0 + w_1 + (x_0 + x_1)i + (y_0 + y_1)j + (z_0 + z_1)k
			\\
			& q_0 - q_1
			= q_0 + (-q_1)
			= (w_0 + x_0i + y_0j + z_0k) + (-w_1 - x_1i - y_1j - z_1k)
			\\
			& \hspace{3em} = w_0 - w_1 + (x_0 - x_1)i + (y_0 - y_1)j + (z_0 - z_1)k
			\\
			& q_0q_1 = (w_0 + x_0i + y_0j + z_0k)(w_1 + x_1i + y_1j + z_1k)
			\\ & \hspace{1.8em} = w_0w_1 + w_0x_1i + w_0y_1j + w_0z_1k + w_1x_0i + x_0ix_1i + x_0iy_1j + x_0iz_1k
			\\ & \hspace{1.8em} + w_1y_0j + y_0jx_1i + y_0jy_1j + y_0jz_1k + w_1z_0k + z_0kx_1i + z_0ky_1j + z_0kz_1k
			\\ & \hspace{1.8em} = w_0w_1 + w_0x_1i + w_0y_1j + w_0z_1k + w_1x_0i - x_0x_1 + x_0y_1k - x_0z_1j + w_1y_0j
			\\ & \hspace{1.8em} - y_0x_1k - y_0y_1 + y_0z_1i + w_1z_0k + z_0x_1j - z_0y_1i - z_0z_1
			\\ & \hspace{1.8em} = w_0w_1 - x_0x_1 - y_0y_1 - z_0z_1 + i(w_0x_1 + w_1x_0 + y_0z_1 - z_0y_1)
			\\ & \hspace{1.8em} + j(w_0y_1 - x_0z_1 + w_1y_0 + z_0x_1) + k(w_0z_1 + x_0y_1 - y_0x_1 + w_1z_0)
		\end{aligned}
	\]
	
	\section{Magnitude}
	
	Just like in complex numbers or vectors, quaternions also have a norm or magnitude that is computed right in the same way or using the same logic as them. By using the pythagorean theorem:
	
	\[
		\|q\|^2 = w^2 + x^2 + y^2 + z^2 \implies |q| = \sqrt{w^2 + x^2 + y^2 + z^2}
	\]
	
	\section{Conjugate of a quaternion}
	
	The conjugate of a quaternion $q$ is the same concept as the conjugate of a complex number, we just flip the symbol from the imaginary components:
	
	\[
		q^* = w - xi - yj - zk
	\]
	
	There is a interesting property of the conjugate, if we multiply $qq^*$, we get:
	
	\[
		\begin{aligned}
			& qq^* = (w + xi + yj + zk)(w - xi - yj - zk)
			\\
			& \hspace{1.4em} = w^2 - wxi - wyj - wzk + wxi + x^2 - xyk + xzj + wyj + yxk + y^2- yzi
			\\
			& \hspace{1.4em} + wzk - zxj + zyi + z^2
			\\
			& \hspace{1.4em} = w^2 + x^2 + y^2 + z^2 = |q|^2
		\end{aligned}
	\]
	
	By expanding all terms we see that the majority of them can be canceled out, returning simply the squared magnitude of the quaternion.
	
	\section{Inverse of a quaternion}
	
	The interesting part here is nothing but the inverse as the inverse is defined as another quaternion $q^{-1}$ such as $q^{-1}q = 1$:
	
	\[
		q^{-1}qq^* = q^* \implies q^{-1}|q|^2 = q^* \implies q^{-1} = \frac{q^*}{|q|^2}
	\]
	
	\[
		\therefore q^{-1} = \frac{q^*}{|q|^2}
	\]
	
	We see that if $|q|^2 = 1$, which implies that $|q| = 1$, $q^{-1} = q^*$.
	
\end{document}